\subsection{Ideals of Simple Rings}

Let D be a field and V a right vector space over D of finite dimension \(n \geqslant 1\) . Consider the simple ring \(A = \text{End}_{\text{D}}(\text{V})\) . For any linear subspace W of V, we denote the set of elements a of A satisfying \(aW = 0\) (resp. \(aV \subset W\) ) by \(\mathfrak{a}(\text{W})\) (resp. \(\mathfrak{b}(\text{W})\) ).

PROPOSITION 7. --- a) The mapping \(W \mapsto \mathfrak{a}(W)\) is a bijection from the set of linear subspaces of V to the set of left ideals of A.

\begin{enumerate}
\def\labelenumi{\alph{enumi})}
\setcounter{enumi}{1}
\item
  The mapping \(W \mapsto \mathfrak{b}(W)\) is a bijection from the set of linear subspaces of V to the set of right ideals of A.
\item
  Let \(\mathrm{W}_1\) and \(\mathrm{W}_2\) be linear subspaces of \(\mathrm{V}\) . The relations \(\mathrm{W}_1 \subset \mathrm{W}_2\) , \(\mathfrak{a}(\mathrm{W}_1) \supset \mathfrak{a}(\mathrm{W}_2)\) , and \(\mathfrak{b}(\mathrm{W}_1) \subset \mathfrak{b}(\mathrm{W}_2)\) are equivalent.
\end{enumerate}

Assertion b) follows from Example 1, b) of VIII, p.~104 applied to the invertible \((\mathrm{D}^{\circ},\mathrm{A}^{\circ})\) -bimodule V, as does the equivalence of the relations \(\mathrm{W}_1\subset \mathrm{W}_2\) and \(\mathfrak{b}(\mathrm{W}_1)\subset \mathfrak{b}(\mathrm{W}_2)\) .

Let \(V^{*}\) be the dual of V, viewed as a right vector space over the opposite field \(D^{\circ}\) of D. For any subspace W of V, denote the orthogonal of W in \(V^{*}\) by \(W'\) . The mapping \(W \mapsto W'\) is a bijection from the set of subspaces of V to the set of subspaces of \(V^{*}\) . If \(W_{1}\) and \(W_{2}\) are two subspaces of V, then the relations \(W_{1} \subset W_{2}\) and \(W_{1}' \supset W_{2}'\) are equivalent. Now, the mapping \(u \mapsto {}^{t}u\) is an isomorphism from A to the opposite ring of \(\operatorname{End}_{D^{\circ}}(V^{*})\) ; it transforms left ideals of A into right ideals of \(\operatorname{End}_{D^{\circ}}(V^{*})\) and \(\mathfrak{a}(W)\) into the set \(\mathfrak{b}(W')\) of endomorphisms h of \(V^{*}\) such that \(h(V^{*}) \subset W'\) . Assertion a), as well as the equivalence of the relations \(W_{1} \subset W_{2}\) and \(\mathfrak{a}(W_{1}) \supset \mathfrak{a}(W_{2})\) , then follows from the assertion analogous to b) for the dual \(V^{*}\) of V.

COROLLARY. --- a) The minimal left ideals of A are the ideals \(\mathfrak{a}(\mathrm{H})\) , where H is a hyperplane in V. The maximal left ideals of A are the ideals \(\mathfrak{a}(\mathrm{L})\) , where L is a line in V.

\begin{enumerate}
\def\labelenumi{\alph{enumi})}
\setcounter{enumi}{1}
\tightlist
\item
  The minimal right ideals of A are the ideals \(\mathfrak{b}(\mathrm{L})\) , where L is a line in V. The maximal right ideals of A are the ideals \(\mathfrak{b}(\mathrm{H})\) , where H is a hyperplane in V.
\end{enumerate}

Let \((\mathrm{L}_{i})_{i\in\mathrm{I}}\) be a family of lines with direct sum V. Let \((\varepsilon_{i})_{i\in\mathrm{I}}\) be the family of projectors associated with the decomposition \(V=\oplus_{i\in\mathrm{I}}L_{i}\) . The \(\varepsilon_{i}\) are idempotent in A, and we have \(\varepsilon_{i}\varepsilon_{j} = 0\) for \(i\neq j\) and \(\sum_{i\in \mathrm{I}}\varepsilon_i = 1\) . Denote the hyperplane \(\sum_{j\neq i}\mathrm{L}_j\) by \(\mathrm{H}_i\) ; it is the kernel of \(\varepsilon_{i}\) . We then have

\[
\mathfrak {a} (\mathrm{H} _ {i}) = \mathrm{A}   \varepsilon_ {i}  , \quad \mathfrak {b} (\mathrm{L} _ {i}) = \varepsilon_ {i}   \mathrm{A}  .
\]

The A-module \(A_{s}\) is the direct sum of the family \((\mathfrak{a}(\mathrm{H}_{i}))_{i\in\mathrm{I}}\) of minimal left ideals, and \(A_{d}\) is the direct sum of the family \((\mathfrak{b}(\mathrm{L}_{i}))_{i\in\mathrm{I}}\) of minimal right ideals.

Consider the specific case \(V = D_{d}^{n}\) , and identify A with the matrix ring \(\mathbf{M}_{n}(\mathrm{D})\) . Denote the interval \([1, n]\) in N by I and the canonical basis of V by \((v_{i})_{i \in \mathrm{I}}\) ; set \(L_{i} = v_{i}D\) , and denote by \(E_{ij}\) the matrix units (II, §10, No.~3, p.~341). We then have \(\varepsilon_{i} = E_{ii}\) . The left ideal \(AE_{ii}\) is equal to \(DE_{1i} + \cdots + DE_{ni}\) and consists of the matrices with all columns except the i-th equal to zero. The right ideal \(E_{ii}A\) is equal to \(DE_{i1} + \cdots + DE_{in}\) and consists of the matrices with all lines except the ith equal to zero. We also have the relation

\[
\mathrm{E} _ {i i} \mathrm{AE} _ {j j} = \mathrm{E} _ {i i} \mathrm{A} \cap \mathrm{AE} _ {j j} = \mathrm{DE} _ {i j}
\]

for i and j between 1 and n.
% semantic-fragments: legacy-input-seam
\relax{}
